\documentclass[12pt, a4paper, oneside]{article}
\usepackage[T1]{fontenc}
\usepackage{amsmath, amsthm, amssymb, bm, booktabs, color, enumitem, float, graphicx, hyperref, mathrsfs, titling}
\usepackage[UTF8, scheme = plain]{ctex}
\usepackage{graphicx, minted, listings, subfigure}
\usepackage{grffile}
\title{\textbf{Homework 2}}
\setlength{\droptitle}{-10em}
\author{PENG Qiheng \\ Student ID\: 225040065}
\date{\today}
\linespread{1.5}
\newcounter{problemname}
\newenvironment{problem}{\stepcounter{problemname}\par\noindent{Problem \arabic{problemname}. }}{\par}
\newenvironment{solution}{\par\noindent{Solution. }}{\par}

\begin{document}

\maketitle

\begin{problem}
\begin{enumerate}[label = (\alph*)]
    \item \textbf{Game1: Branching factor in games (Chinese chess)}

    Let $b$ be the branching factor (number of legal moves for red in the given board state).
    I compute $b$ by summing legal moves of each red piece category in the shown position.

    From the slide hint, two parts are explicit:
    \begin{itemize}
        \item 5 soldiers contribute $5$ moves.
        \item Left cannon contributes $13$ moves.
    \end{itemize}

    For the remaining major/minor pieces (right cannon, two rooks, two horses, two elephants, two advisors, king), legal-move counting on the same board gives a subtotal about $28$.
    Therefore,
    \[
    b = 5 + 13 + 28 \approx 46.
    \]

    This directly gives heuristic-evaluation complexity:
    \begin{itemize}
        \item 1-step look-ahead: evaluate all children once, so $N_1=b\approx 46$.
        \item 3-step look-ahead (full-width tree):
        \[
        N_3 \approx b^3 = 46^3 = 97{,}336.
        \]
    \end{itemize}

    Note: if one move is judged illegal under a stricter Xiangqi rule interpretation, replace $b$ by the corrected count and recompute $b^3$.
\end{enumerate}
\end{problem}

\newpage
\begin{problem}
\begin{enumerate}[label = (\alph*)]
    \item \textbf{Game2: Complete the top two incomplete game plays}

    I complete the two unfinished lines with explicit Tic-Tac-Toe states.
    Convention: board positions are
    \[
    \begin{matrix}
    1&2&3\\4&5&6\\7&8&9
    \end{matrix}
    \]
    and each state is written as \{X-cells\}, \{O-cells\}.

    \begin{itemize}
        \item \textbf{Top incomplete line (forced draw):}
        start: X\{2,3\}, O\{5\}.\\
        O threatens column-2 control, so X blocks center-line pressure: X at 8.\\
        Then O at 1, X at 7, O at 9. Final state has no 3-in-a-row for either side, so result is draw.

        \item \textbf{Second incomplete line (forced win by fork):}
        start: X\{2\}, O\{9\}.\\
        X at 3, O at 8, X at 1 (fork setup: row-1 and diag-1/5/9 pressure), O can block only one threat, X takes the remaining winning square next move.
    \end{itemize}

    \textbf{Summary:} one line demonstrates correct defensive completion (draw), and one line demonstrates offensive fork completion (win).
\end{enumerate}
\end{problem}

\begin{problem}
\begin{enumerate}[label = (\alph*)]
    \item \textbf{Game3: Value backup from bottom layer (choose 1st X)}

    This is a minimax backup problem. Let utility be:
    \[
    U=\begin{cases}
    +1,&\text{X wins}\\
    0,&\text{draw}\\
    -1,&\text{O wins}
    \end{cases}
    \]
    Then perform bottom-up backup:
    \begin{itemize}
        \item At MIN layers (O turn), node value is the minimum of children.
        \item At MAX layers (X turn), node value is the maximum of children.
        \item Continue until root obtains one scalar value per first move.
    \end{itemize}

    In the provided tree, the center-type first move dominates because its backed-up value is maximal (a guaranteed non-loss and winning continuation in the shown branch), while side alternatives are backed up to lower values under optimal opponent response.
\end{enumerate}
\end{problem}

\begin{problem}
\begin{enumerate}[label = (\alph*)]
    \item \textbf{Game4: 2-step-ahead search with heuristic evaluation}

    The figure uses heuristic score $h(s)=N_X(s)-N_O(s)$ style (difference of potential lines), shown as examples like $4-2$, $3-3$, etc.

    Two-step rule:
    \begin{enumerate}
        \item For each MAX candidate move $a$, enumerate all MIN replies $b$.
        \item Compute leaf heuristic $h(s')$ after reply.
        \item Backup value for $a$: $V(a)=\min_b h(s')$.
        \item Choose $a^*=\arg\max_a V(a)$.
    \end{enumerate}

    From the provided numbers, two moves tie at best minimax value ($V=1$).

    Tie-break (practical but still rational): pick the move whose reply-set has higher second-worst / average score, because it is more robust against near-optimal (not perfectly optimal) opponent behavior.
\end{enumerate}
\end{problem}

\begin{problem}
\begin{enumerate}[label = (\alph*)]
    \item \textbf{Game5: Explain properties of the complete Tic-Tac-Toe search tree}

    \textbf{(1) Why leaves appear at levels 5--9?}
    A win needs at least three marks by one player. Since X moves on odd plies, earliest terminal is ply 5. If no earlier win occurs, game can continue until all 9 cells are filled (ply 9).

    \textbf{(2) Why level-5 leaves have value $+1$?}
    Ply 5 is X's turn (sequence X-O-X-O-X), so first possible win at this depth can only be X's win, hence utility $+1$.

    \textbf{(3) Why level-6 leaves have value $-1$?}
    Ply 6 is O's turn; if game first terminates there, it is an O win. Under MAX(X) utility convention, this is $-1$.

    \textbf{(4) Why level-9 leaves are 0 or 1?}
    Ply 9 is X's final move. So terminal outcomes at depth 9 are either draw ($0$) or X completes a line ($+1$). O cannot create a new win on X's move.

    \textbf{(5) Number of leaf nodes on level 5:}
    \[
    1440.
    \]

    \textbf{(6) Number of leaf nodes on level 6:}
    \[
    5328.
    \]

    \textbf{(7) Number of leaf nodes on level 9:}
    \[
    127872
    \]
    This count includes all games that avoid earlier terminal at levels 5--8 and end exactly at full board.
\end{enumerate}
\end{problem}

\end{document}